\documentclass[11pt]{article}

\usepackage[margin=1in]{geometry}
\usepackage{amsmath,amssymb,bm}
\usepackage{physics}
\usepackage{hyperref}

\newcommand{\Jzz}{J_{zz}}
\newcommand{\Jpm}{J_{\pm}}
\newcommand{\Jpmpm}{J_{\pm\pm}}
\newcommand{\Jthree}{J_{3s}}
\newcommand{\rr}{\mathbf r}
\newcommand{\ee}{\mathbf e}
\newcommand{\sgn}{\operatorname{sgn}}

\begin{document}

\title{Local-frame non-Kramers pyrochlore Hamiltonian and symmetry covariance}
\author{}
\date{}
\maketitle

\section{Local non-Kramers pseudospins}

The rare-earth sites of the pyrochlore lattice are labelled by a Bravais cell
$R$ and a sublattice index $\mu=0,1,2,3$.  Each sublattice carries its own
local orthonormal frame
\begin{equation}
  \left(\hat{\bm x}_{\mu},\hat{\bm y}_{\mu},\hat{\bm z}_{\mu}\right),
  \qquad
  \hat{\bm z}_{\mu}\parallel \langle 111\rangle ,
\end{equation}
with the four local Ising axes chosen as
\begin{equation}
\begin{aligned}
  \hat{\bm z}_0 &= \frac{(1,1,1)}{\sqrt 3},&
  \hat{\bm z}_1 &= \frac{(1,-1,-1)}{\sqrt 3},\\
  \hat{\bm z}_2 &= \frac{(-1,1,-1)}{\sqrt 3},&
  \hat{\bm z}_3 &= \frac{(-1,-1,1)}{\sqrt 3}.
\end{aligned}
\end{equation}
The low-energy doublet is represented by pseudospin-$1/2$ operators
$S^\alpha_{R\mu}$ in this local frame.  For a non-Kramers doublet,
$S^z$ transforms as the magnetic dipole component along the local Ising axis,
whereas $S^x$ and $S^y$ transform as electric quadrupoles.  Equivalently,
\begin{equation}
  S^\pm_{R\mu}=S^x_{R\mu}\pm iS^y_{R\mu}
\end{equation}
are local quadrupolar raising/lowering operators.

The key point for symmetry is that a space-group operation $g$ does not act as
a pure site permutation on the pseudospin components.  It maps
$(R,\mu)\mapsto (R',\mu')$ and rotates the local transverse frame by an angle
$\chi_g(\mu)$ about the local $\hat{\bm z}$ axis.  In the local basis one may
write
\begin{equation}
\begin{aligned}
  g: S^z_{R\mu} &\mapsto S^z_{R'\mu'},\\
  g: S^+_{R\mu} &\mapsto e^{+i\chi_g(\mu)} S^+_{R'\mu'},\\
  g: S^-_{R\mu} &\mapsto e^{-i\chi_g(\mu)} S^-_{R'\mu'}.
\end{aligned}
\label{eq:local-frame-transform}
\end{equation}
The sublattice-dependent phases are the local-frame expression of the
standard non-Kramers pyrochlore symmetry action.  A Hamiltonian written in
local-frame operators is invariant when its bond and multi-spin phase factors
compensate these local-frame rotations.

\section{Nearest-neighbor non-Kramers XYZ Hamiltonian}

The full nearest-neighbor non-Kramers pyrochlore Hamiltonian in the local
pseudospin basis can be written as
\begin{equation}
  H_{\rm XYZ}
  =
  \sum_{\langle i j\rangle}
  \biggl[
  \Jzz S^z_i S^z_j
  -
  \Jpm\left(S^+_iS^-_j+S^-_iS^+_j\right)
  +
  \Jpmpm\left(
    \gamma_{ij} S^+_iS^+_j
    +
    \gamma_{ij}^{*} S^-_iS^-_j
  \right)
  \biggr].
  \label{eq:hxyz}
\end{equation}
Here $i\equiv(R,\mu)$ and $j\equiv(R',\nu)$, and $\gamma_{ij}$ is the standard
bond-dependent non-Kramers phase factor.  In a common sublattice convention
$\gamma_{ij}\in\{1,\omega,\omega^2\}$ with $\omega=e^{2\pi i/3}$; the precise
assignment depends on the oriented nearest-neighbor bond.  This
$\gamma_{ij}$ factor is required because $S^\pm$ are local-frame quadrupolar
operators.  Under Eq.~\eqref{eq:local-frame-transform},
$S^+_iS^+_j$ acquires the phase
$e^{i[\chi_g(\mu_i)+\chi_g(\mu_j)]}$, and the transformed bond coefficient
$\gamma_{g i,g j}$ compensates this phase so that
\begin{equation}
  \gamma_{ij} S_i^+S_j^+
  \mapsto
  \gamma_{g i,g j} S_{g i}^+S_{g j}^+ .
\end{equation}
The Hermitian-conjugate $S^-S^-$ term transforms with the complex conjugate
factor.  Thus the complete sum in Eq.~\eqref{eq:hxyz} is a scalar under the
pyrochlore space group.

This form is often called the local-frame XYZ model.  In Cartesian local
coordinates it is equivalent to bond-dependent exchange in the transverse
$xy$ plane.  The relation to local Cartesian couplings is, for the convention
used by the present helper,
\begin{equation}
  \Jpm = -\frac{J_{xx}+J_{yy}}{4}.
\end{equation}
When $J_{xx}\neq J_{yy}$, the transverse anisotropy appears as a nonzero
$\Jpmpm$ channel with the sublattice/bond phase $\gamma_{ij}$.

\subsection{XXZ limit used in the present scans}

The numerical scans discussed here use the XXZ submanifold
\begin{equation}
  \Jpmpm=0,\qquad
  H_{\rm XXZ}
  =
  \sum_{\langle i j\rangle}
  \left[
  \Jzz S^z_i S^z_j
  -
  \Jpm\left(S^+_iS^-_j+S^-_iS^+_j\right)
  \right].
  \label{eq:hxxz}
\end{equation}
For the $jpm03$ point, $J_{xx}=J_{yy}=0.6\Jzz$, hence $\Jpm=-0.3\Jzz$ and
$\Jpmpm=0$.

For ordinary periodic boundary conditions, all boundary twist phases are set
to zero.  More generally, the finite-cluster twist implementation multiplies
the transverse hopping on a wrapped bond by the Peierls phase
\begin{equation}
  S^+_iS^-_j \mapsto
  e^{-i\bm n_{ij}\cdot\bm\varphi}S^+_iS^-_j,
\end{equation}
where $\bm n_{ij}$ is the integer wrap vector of the minimum-image bond.
For the normal PBC calculations discussed here, $\bm\varphi=\bm 0$.

In the local-frame convention, Eq.~\eqref{eq:hxxz} is the $\Jpmpm=0$ scalar
limit of Eq.~\eqref{eq:hxyz}.  The full codebase contains non-Kramers
$\Jpmpm$ phase-factor logic in the general helper path, while
\texttt{twist\_trilinear\_helper.py} deliberately writes the XXZ two-body
limit plus the Kadowaki three-spin term used below.

\section{Kadowaki three-spin term}

The trilinear perturbation is the Kadowaki three-spin interaction
\begin{equation}
  H_{3s}
  =
  \sum_{a=1}^{3} J_{3s,a}
  \sum_{\langle r,r',r''\rangle_a}
  \left[
    e^{i\phi^{(a)}_{r,r'r''}}
    S^+_{r}S^z_{r'}S^z_{r''}
    +
    e^{-i\phi^{(a)}_{r,r'r''}}
    S^-_{r}S^z_{r'}S^z_{r''}
  \right].
  \label{eq:h3s}
\end{equation}
The index $a=1,2,3$ labels the three Kadowaki geometric classes of unordered
nearest-neighbor pairs around the central site $r$:
\begin{itemize}
  \item $a=1$: collinear pair, where $r'$ and $r''$ have the same sublattice
        index and lie on opposite tetrahedra sharing $r$;
  \item $a=2$: same-tetrahedron pair, where $r'$ and $r''$ are both neighbors
        of $r$ in the same tetrahedron;
  \item $a=3$: cross pair, where $r'$ and $r''$ lie on opposite tetrahedra
        adjacent to $r$.
\end{itemize}
For each central site there are $3+6+6=15$ unordered pairs.  The code supports
either independent $(J_{3s,1},J_{3s,2},J_{3s,3})$ or a scalar $J_3$ broadcast
to all three classes.

The phases are roots of unity,
\begin{equation}
  e^{i\phi^{(a)}}=\omega^{p^{(a)}},
  \qquad
  \omega=e^{2\pi i/3},
\end{equation}
with the integers $p^{(a)}$ read from the Kadowaki phase tables in the
sublattice convention $\mu=0,1,2,3$.  The implementation uses
\begin{equation}
\begin{array}{c|ccc}
\mu_{\rm center} & \text{outer sublattice }1 & 2 & 3 \\ \hline
0 & -1 & +1 & 0
\end{array}
\end{equation}
and the corresponding cyclic tables for the other central sublattices for the
collinear class, together with the pair tables for the same-tetrahedron and
cross classes.  These phase tables are exactly what makes
Eq.~\eqref{eq:h3s} a local-frame scalar.  Under a space-group operation, the
central transverse operator transforms as
$S^+_r\mapsto e^{i\chi_g}S^+_{g r}$, while the Kadowaki phase transforms as
\begin{equation}
  e^{i\phi^{(a)}_{r,r'r''}}
  \mapsto
  e^{-i\chi_g(r)} e^{i\phi^{(a)}_{g r,g r',g r''}},
\end{equation}
so that the product
$e^{i\phi^{(a)}}S^+S^zS^z$ maps into the corresponding term in the same
Hamiltonian.  The Hermitian-conjugate $S^-$ row transforms with the complex
conjugate phase.  Thus the complete sum over all triplets in
Eq.~\eqref{eq:h3s} is invariant under the pyrochlore space group.

The present code writes the file order
\begin{equation}
  S^z_{r'}\,S^\pm_{r}\,S^z_{r''},
\end{equation}
with the same complex coefficient as Eq.~\eqref{eq:h3s}.  Since the three
operators act on distinct sites, this is the same operator product.

\section{Quadrupolar strain probes}

The thermal-expansion proxy is computed for five local-XY quadrupole
operators:
\begin{equation}
  Q_1,\quad Q_2,\quad Q_{xy},\quad Q_{xz},\quad Q_{yz}.
\end{equation}
They are written as
\begin{equation}
  Q = \sum_{R,\mu}
  \left[
    c^x_{\mu}(Q)\,S^x_{R\mu}
    +
    c^y_{\mu}(Q)\,S^y_{R\mu}
  \right],
  \label{eq:q-general}
\end{equation}
with sublattice-dependent coefficients.  The convention used in the operator
files is
\begin{equation}
\begin{array}{c|rrrr|rrrr}
  & \multicolumn{4}{c|}{c^x_\mu}
  & \multicolumn{4}{c}{c^y_\mu} \\
  Q & \mu=0 & 1 & 2 & 3 & \mu=0 & 1 & 2 & 3 \\ \hline
  Q_1
    & \sqrt3 & \sqrt3 & \sqrt3 & \sqrt3
    & -1 & -1 & -1 & -1 \\
  Q_2
    & -1 & -1 & -1 & -1
    & -\sqrt3 & -\sqrt3 & -\sqrt3 & -\sqrt3 \\
  Q_{xy}
    & +\frac14 & -\frac14 & -\frac14 & +\frac14
    & +\frac{\sqrt3}{4} & -\frac{\sqrt3}{4}
    & -\frac{\sqrt3}{4} & +\frac{\sqrt3}{4} \\
  Q_{xz}
    & +\frac14 & -\frac14 & +\frac14 & -\frac14
    & -\frac{\sqrt3}{4} & +\frac{\sqrt3}{4}
    & -\frac{\sqrt3}{4} & +\frac{\sqrt3}{4} \\
  Q_{yz}
    & -\frac12 & -\frac12 & +\frac12 & +\frac12
    & 0 & 0 & 0 & 0
\end{array}
\label{eq:q-coeff-table}
\end{equation}
Thus, explicitly,
\begin{equation}
\begin{aligned}
  Q_1
  &= \sum_{R,\mu}
     \left(\sqrt3\,S^x_{R\mu}-S^y_{R\mu}\right),\\
  Q_2
  &= -\sum_{R,\mu}
     \left(S^x_{R\mu}+\sqrt3\,S^y_{R\mu}\right),
\end{aligned}
\label{eq:eg-explicit}
\end{equation}
and
\begin{equation}
\begin{aligned}
  Q_{xy}
  &= \sum_R \biggl[
      \frac14 S^x_{R0}-\frac14 S^x_{R1}
      -\frac14 S^x_{R2}+\frac14 S^x_{R3}
      +\frac{\sqrt3}{4} S^y_{R0}
      -\frac{\sqrt3}{4} S^y_{R1}
      -\frac{\sqrt3}{4} S^y_{R2}
      +\frac{\sqrt3}{4} S^y_{R3}
     \biggr],\\
  Q_{xz}
  &= \sum_R \biggl[
      \frac14 S^x_{R0}-\frac14 S^x_{R1}
      +\frac14 S^x_{R2}-\frac14 S^x_{R3}
      -\frac{\sqrt3}{4} S^y_{R0}
      +\frac{\sqrt3}{4} S^y_{R1}
      -\frac{\sqrt3}{4} S^y_{R2}
      +\frac{\sqrt3}{4} S^y_{R3}
     \biggr],\\
  Q_{yz}
  &= \sum_R \left[
      -\frac12 S^x_{R0}
      -\frac12 S^x_{R1}
      +\frac12 S^x_{R2}
      +\frac12 S^x_{R3}
     \right].
\end{aligned}
\label{eq:t2g-explicit}
\end{equation}
The signs above include the strain-operator convention used in
\texttt{make\_quad\_operator\_files.py}: $O_\mu=+Q_1$, $O_\nu=-Q_2$, and
$O_{ij}=-\mathcal Q_{ij}$ relative to the manuscript quadrupoles.

For the legacy ED operator files, Eq.~\eqref{eq:q-general} is converted to
raising/lowering rows through
\begin{equation}
  c^x S^x + c^y S^y
  =
  \left(\frac{c^x}{2}-\frac{i c^y}{2}\right)S^+
  +
  \left(\frac{c^x}{2}+\frac{i c^y}{2}\right)S^- .
\end{equation}
The two $E_g$ operators are uniform local transverse sums, while the three
$T_{2g}$ operators are sublattice-staggered local-XY combinations
corresponding to $Q_{xy}$, $Q_{xz}$, and $Q_{yz}$ shear channels.  These
operators transform as irreducible components of cubic strain.  Consequently,
a cubic-symmetric Hamiltonian cannot mix the $E_g$ and $T_{2g}$ response
blocks.

The quantity plotted as a thermal-expansion coefficient is the connected
energy--quadrupole covariance,
\begin{equation}
  \alpha_Q(T)
  =
  \frac{\langle QH\rangle_T-\langle Q\rangle_T\langle H\rangle_T}{T^2}.
  \label{eq:alphaq}
\end{equation}
With the usual canonical derivative convention,
\begin{equation}
  \frac{d\langle Q\rangle_T}{dT}
  =
  \frac{\langle QH\rangle_T-\langle Q\rangle_T\langle H\rangle_T}{T^2},
\end{equation}
so Eq.~\eqref{eq:alphaq} is the temperature derivative of the thermal
quadrupole expectation value in the same sign convention as the Hamiltonian.

\section{Implementation status and caveat}

For ordinary, non-symmetry-projected ED/FTLM runs, the generated Hamiltonian
files used in the present scans implement the XXZ limit
Eq.~\eqref{eq:hxxz} and the Kadowaki term Eq.~\eqref{eq:h3s}.  In particular:
\begin{itemize}
  \item the PBC two-body file \texttt{InterAll.dat} contains the $S^zS^z$ and
        $S^+S^-+\mathrm{h.c.}$ XXZ terms;
  \item the trilinear file \texttt{ThreeBodyG.dat} contains $480$ rows for the
        $16$-site PBC cluster when $J_3\neq0$, corresponding to
        $16\times15$ triplets and their Hermitian conjugates;
  \item the ordinary QED path prints
        \texttt{Loading three-body terms ... ThreeBodyG.dat}, confirming that
        the Kadowaki term is included in non-\texttt{--symm} calculations.
\end{itemize}

There is a separate computational caveat: the current QED streaming
\texttt{--symm} path should not be used for the three-body Kadowaki
Hamiltonian until \texttt{ThreeBodyG.dat} is both loaded and included in the
automorphism validation.  This caveat does not change the physical symmetry
argument above; it only restricts which numerical execution path is safe.

\section*{Literature conventions}

The local-frame non-Kramers pseudospin convention follows the standard
pyrochlore treatments of Onoda--Tanaka, Savary--Balents, and related
non-Kramers quantum spin ice literature, where $S^z$ is dipolar and
$(S^x,S^y)$ are quadrupolar components.  The three-spin interaction and its
phase tables follow Kadowaki \emph{et al.}, Phys. Rev. B \textbf{105},
014439 (2022).

\end{document}
